\chapter{Desarrollo Teórico}

\section{Índice Horario}

El índice horario es una convención utilizada para indicar el 
desfase existente entre la fuerza electromotriz (FEM) primaria y la secundaria 
de un transformador trifásico. Este desfase depende del sentido de arrollamiento 
de las bobinas y de la forma en que estas se conectan entre sí.

La nomenclatura se compone de:
\begin{itemize}
    \item Una letra mayúscula que indica la conexión del primario ($Y$: estrella, 
    $D$: triángulo, $Z$: zigzag).
    \item Una letra minúscula que indica la conexión del secundario ($y$: estrella, 
    $d$: triángulo, $z$: zigzag).
    \item Un número (0 al 11) que multiplicado por $30^\circ$ nos da el ángulo de 
    retraso del vector de baja tensión respecto al de alta tensión.
\end{itemize}

\section{Representación del Reloj}

Se utiliza la analogía de las agujas de un reloj para visualizar el desfase. El 
vector de tensión del bobinado de alta tensión se sitúa siempre a las 12 (posición 
cero). El vector de tensión del bobinado de baja tensión se sitúa en la hora que 
indica el índice.

Cada hora representa un desfase de $30^\circ$ en sentido horario (retraso).

\begin{figure}[!ht]
    \centering
    \begin{tikzpicture}[line cap=rect, line join=round, >=Stealth]
    % Clock face
    \draw[thick] (0,0) circle (2cm);
    \foreach \x in {1,2,...,12} {
        \draw ({90-\x*30}:1.85cm) -- ({90-\x*30}:2cm);
        \node at ({90-\x*30}:1.6cm) {\footnotesize \x};
    }
    
    % AT Phasor (fixed at 12) - Longer to stand out
    \draw[->, ultra thick, blue] (0,0) -- (0,2.6) node[above, font=\small\bfseries] {AT ($U_A$)};
    
    % BT Phasor (at 11)
    \draw[->, ultra thick, red] (0,0) -- (120:2.3) node[above, font=\small\bfseries] {BT ($u_a$)};
    
    % Angle arc (from 12 o'clock clockwise to 11 o'clock)
    \draw[->, thick] (0,0.8) arc (90:-240:0.8);
    
    % Angle label - Positioned in a clear area (around 5 o'clock)
    \node[draw=black!20, fill=white, rounded corners=2pt, font=\footnotesize, inner sep=2pt] 
        at (315:1.2) {$330^\circ = 11 \times 30^\circ$};
    
\end{tikzpicture}

    \caption{Representación del índice horario mediante el reloj.}
    \label{fig:clock}
\end{figure}

\section{Conexiones Comunes}

Las combinaciones de conexiones permiten obtener diferentes índices horarios. 
Los más utilizados en la industria son:

\begin{itemize}
    \item \textbf{Yy0:} Ambas conexiones en estrella. No hay desfase. Se utiliza en 
    redes de transporte pero presenta problemas con la tercera armónica si no 
    tiene devanado de compensación.
    \item \textbf{Dy11:} Primario en triángulo, secundario en estrella. El 
    secundario adelanta $30^\circ$ al primario (o retrasa $330^\circ$). Es el más 
    común en distribución, permitiendo neutro accesible y equilibrando cargas.
    \item \textbf{Yz5:} Primario en estrella, secundario en zigzag. Permite 
    un fuerte desequilibrio de cargas en el secundario.
\end{itemize}

\section{Resumen de Grupos y Aplicaciones}

A continuación se presenta una tabla resumen con los grupos de conexión más 
frecuentes y sus aplicaciones principales:

\begin{table}[!ht]
    \centering
    \begin{tabular}{|l|l|l|p{6cm}|}
        \hline
        \textbf{Grupo} & \textbf{Desfase} & \textbf{Conexión} & \textbf{Aplicación Típica} \\ \hline
        Yy0 & $0^\circ$ & Estrella-Estrella & Grandes sistemas de transmisión, requiere devanado terciario para armónicos. \\ \hline
        Dd0 & $0^\circ$ & Triángulo-Triángulo & Transformadores de gran corriente y baja tensión. \\ \hline
        Dy11 & $330^\circ$ & Triángulo-Estrella & Distribución urbana e industrial (el más común). Permite neutro. \\ \hline
        Yz5 & $150^\circ$ & Estrella-Zigzag & Distribución con cargas muy desequilibradas. \\ \hline
        Dz0 & $0^\circ$ & Triángulo-Zigzag & Aplicaciones especiales donde se requiere mitigar armónicos y desequilibrios. \\ \hline
    \end{tabular}
    \caption{Grupos de conexión y aplicaciones.}
    \label{tab:indices}
\end{table}

\section{Ventajas y Aplicaciones}

La elección del índice horario es determinante para diversas aplicaciones:

\begin{enumerate}
    \item \textbf{Puesta en paralelo:} Para conectar dos transformadores en 
    paralelo, es condición necesaria que ambos pertenezcan al mismo grupo de 
    índice horario (o grupos compatibles mediante cambio de bornes). De lo 
    contrario, circularían corrientes de cortocircuito inadmisibles.
    \item \textbf{Tratamiento de Armónicos:} Conexiones en triángulo ($D$) 
    permiten la circulación de la tercera armónica dentro de la malla, evitando 
    su propagación a la red.
    \item \textbf{Acceso al Neutro:} La conexión en estrella ($y$) permite 
    conectar el neutro a tierra, facilitando la protección y el suministro a 
    cargas monofásicas.
\end{enumerate}

\section{Puesta en Paralelo}

Para que dos transformadores trifásicos puedan trabajar en paralelo, deben 
cumplir:
\begin{itemize}
    \item Igual relación de transformación (tensiones iguales).
    \item Igual tensión de cortocircuito (para un reparto equitativo de carga).
    \item \textbf{Igual índice horario} (o desfase nulo entre secundarios).
    \item Misma secuencia de fases.
\end{itemize}
